\section{Conclusion}

We began by observing a striking paradox: LLMs that pass 99\% of benchmark tests still fail to fix 60\% of their own compiler errors when given raw error messages. This gap between generation and repair capability motivated our investigation into error message format as an unexplored dimension of self-repair optimization.

\subsection{Summary}

Our work provides the first causal evidence that \emph{how} error information is organized matters independently of \emph{what} information is present. Through controlled experiments isolating format structure from information content, we demonstrated:

\begin{enumerate}
    \item \textbf{Representational alignment works}: Structured error formatting significantly outperforms scrambled formatting with identical content (+14.4\%, $p < 0.001$), proving that organizational structure drives improvement. This finding shifts research focus from \emph{whether} to include feedback to \emph{how} to present it.

    \item \textbf{Scaffolding theory applies}: Intermediate-specificity fix hints (Level 2) achieve optimal repair success (60.2\%), outperforming both no hints (34.7\%) and exact fixes (39.6\%). This inverted-U pattern extends HCI scaffolding principles to LLM guidance design.

    \item \textbf{Simple intervention, substantial impact}: A preprocessing step---reformatting compiler errors with section labels---improves self-repair without model changes, providing a practical path to better code repair systems.
\end{enumerate}

\subsection{Future Directions}

Our findings open several research directions grounded in experimental evidence:

\textbf{Validating pending mechanisms}: The fix specificity inverted-U pattern was confirmed in simulation; real-model validation with resolved CUDA compatibility will strengthen this finding. Similarly, increasing GPT-4 failure samples would enable proper power for scale interaction detection.

\textbf{Extending the format framework}: Our structured format handles 8 Python error types. Extending to statically-typed languages (Java, Rust) and production codebases (SWE-bench) would test generalization. Auto-selection of optimal format per error type is a natural next step.

\textbf{Training-time integration}: While our approach operates at inference time, incorporating structured formats during training could yield compounding benefits. Language-agnostic error representations may enable cross-lingual transfer.

\subsection{Closing Remarks}

The 60\% self-repair failure rate that motivated this work can be significantly reduced by simply reformatting error messages---no model architecture changes, no additional training, no human annotation required. This finding suggests a broader principle: as we integrate external tools with language models, the \emph{interface design} between tool output and model input deserves the same attention we give to model architecture and training data. How we present information to models matters as much as what information we provide.
